\begin{definition}[iid Model]
	In the iid model the algorithm always knows a set of offline vertices $U$, but there is no determined set of online vertices $V$. Instead, there is a set $V'$ of vertex 'archetypes' consisting of a neighborhood $N\subseteq U$ and a probability $p$. When a new vertex $v$ arrives, for each $(N,p)\in V'$ $v$ has of $p$ to have neighborhood $N$. The archetype set $V'$ is known to the algorithm. The amount of vertices that will arrive in total, $n$, is also known. The expected amount of arriving vertices, $p_vn$, is an integer.
\end{definition}
We note two things here:
\begin{enumerate}
	\item obviously, $\sum\limits_{(N,p)\in V'}p=1$
	\item This model is still a game, but the adversary has been stripped of almost all agency and now cannot even control the final graph
\end{enumerate}

Since this is a random graph, the previously defined competitiveness isn't sufficient. We will instead follow \cite{obmStochastic} and use the "approximation factor $\alpha$ with high probability" defined as

\begin{definition}
	For any iid instance $\Gamma$ where $G$ is drawn from it, with high probability $\frac{\mathrm{ALG}(G)}{\mathrm{OPT}(G)}\geq\alpha$
\end{definition}

Which we will refer to as the approximation factor. We note that an approximation ratio of $\alpha$ implies that expected competitiveness $\frac{\E[\mathrm{ALG}]}{\E[\mathrm{OPT}]}$ is at least $\alpha-o(1)$

We will present the proof of two theorems in this chapter:

\begin{theorem}
	No algorithm has an expected competitiveness of $1$ in the iid model, ie there is a constant $c<1$ such that $\alpha<c$ for all algorithms
\end{theorem}
meaning that there is no perfect algorithm, and
\begin{theorem}
	\label{thm:sa-alg}
	The stochastic arrival model is strictly easier than the Semi-Online model since there is an algorithm, TSM, with an approximation factor of $\approx 0.670$
\end{theorem}
In actuality, there is an algorithm with a competitiveness of $0.729$ by \cite{obmStochasticBest}, better than the proven competitiveness of Ranking in the random arrival model. However, this algorithm is very complicated so we will instead showcase a simpler one to show that the $1-\frac{1}{e}$ barrier can be beaten

\section{Competitiveness Limit}
In this section we will follow a proof from \cite{obmStochastic}. Consider a 6-Cycle defined as follows:
\begin{definition}
	$U=\{1,2,3\},V'=\{x=(\set{1,2},\frac{1}{3}),y=(\set{1,3},\frac{1}{3}),z=(\set{2,3},\frac{1}{3})\}$
\end{definition}
\begin{figure}[H]
	\centering
	\begin{tikzpicture}
	\begin{pgfonlayer}{nodelayer}
		\node [style=default] (0) at (-4, 8) {};
		\node [style=default] (1) at (-4, 6) {};
		\node [style=default] (2) at (-4, 4) {};
		\node [style=new] (3) at (0, 8) {};
		\node [style=new] (4) at (0, 6.25) {};
		\node [style=new] (5) at (0, 4) {};
	\end{pgfonlayer}
	\begin{pgfonlayer}{edgelayer}
		\draw (3) to (0);
		\draw (3) to (1);
		\draw (4) to (1);
		\draw (4) to (2);
		\draw (5) to (2);
		\draw (0) to (5);
	\end{pgfonlayer}
\end{tikzpicture}

\end{figure}
If $x$ arrives first, if it is matched to $1$ and then two $y$ vertices $y_1,y_2$ arrive, the second will be unmatchable. In this scenario there is a perfect matching $M^\star=\set{(2,x),(1,y_1),(3,y_2)}$

This argument is symmetric, for any first arrival $v$ matched to $u$ there is an archetype $v'$ that forces this scenario.\\
The probability of this scenario is therefore the probability of $(v',v')$ arriving after $v$, which is $\frac{1}{9}$. Assuming that in all other scenarios the algorithm find the perfect matching, competitiveness is ar most $\frac{1}{9}\frac{2}{3}+\frac{8}{9}1=\frac{26}{27}$.

Therefore, the theorem is proven and the competitiveness of any algorithm in the iid model is bounded by $\approx0.962<1$

\begin{remark}
	We can extend this result asymptotically (albeit not the exact upper bound) by duplicating the 6-Cycle $k$ times
\end{remark}

\section{TSM Algorithm}
In this section we present the $0.729$-competitive algorithm and proof of it's competitiveness from \cite{obmSAlg}